
\documentclass{article}
\usepackage{colortbl}
\usepackage{makecell}
\usepackage{multirow}
\usepackage{supertabular}

\begin{document}

\newcounter{utterance}

\centering \large Interaction Transcript for game `dond', experiment `semi\_en', episode 14 with reflexion{-}Qwen3.5{-}27B.
\vspace{24pt}

{ \footnotesize  \setcounter{utterance}{1}
\setlength{\tabcolsep}{0pt}
\begin{supertabular}{c@{$\;$}|p{.15\linewidth}@{}p{.15\linewidth}p{.15\linewidth}p{.15\linewidth}p{.15\linewidth}p{.15\linewidth}}
    \# & $\;$A & \multicolumn{4}{c}{Game Master} & $\;\:$B\\
    \hline

    \theutterance \stepcounter{utterance}  
    & & \multicolumn{4}{p{0.6\linewidth}}{
        \cellcolor[rgb]{0.9,0.9,0.9}{
            \makecell[{{p{\linewidth}}}]{
                \texttt{\tiny{[P1$\langle$GM]}}
                \texttt{You are playing a negotiation game in which you have to agree on how to divide a set of items among you and another player.} \\
\\ 
\texttt{Rules:} \\
\texttt{(a) You and the other player are given a set of items. Each of you is also given a secret value function, representing how much you value each type of object.} \\
\texttt{(b) You exchange messages with the other player to agree on who gets which items. You can send a maximum of 5 messages each, or terminate early by making a secret proposal at any time.} \\
\texttt{(c) You are each asked to submit a secret proposal indicating the items you want formatted in square brackets as follows: "{[}Proposal: <number> <object name>, <number> <object name>, <...>{]}"} \\
\texttt{(d) If your proposals are complementary, i.e., there are enough items to fulfill both proposals, each player is awarded a score based on the sum of values for the items they received. Otherwise, both of you get zero points.} \\
\texttt{(e) Your goal is to maximize the score you receive.} \\
\\ 
\texttt{Let us start.} \\
\\ 
\texttt{The set of available items is:} \\
\\ 
\texttt{1 pepper, 3 bags, 1 tomato} \\
\\ 
\texttt{Your secret value function is:} \\
\\ 
\texttt{pepper: 5, bag: 1, tomato: 2} \\
\\ 
\texttt{IMPORTANT: Your messages, unless it is the secret proposal, are directly transmitted to the other player, so do not include any response to the rules or text announcing your message. To make a secret proposal, use the indicated format. Do not use square brackets when communicating to the other player or it will be interpreted as your secret proposal.} \\
            }
        }
    }
    & \\ \\

    \theutterance \stepcounter{utterance}  
    & \multicolumn{4}{p{0.6\linewidth}}{
        \cellcolor[rgb]{0.9,0.9,0.9}{
            \makecell[{{p{\linewidth}}}]{
                \texttt{\tiny{[P1$\rangle$GM]}}
                \texttt{I value the pepper the most, followed by the tomato, and the bags are worth the least to me. How do you value these items?} \\
            }
        }
    }
    & & \\ \\

    \theutterance \stepcounter{utterance}  
    & & & \multicolumn{2}{p{0.3\linewidth}}{
        \cellcolor[rgb]{0.9,0.9,0.9}{
            \makecell[{{p{\linewidth}}}]{
                \texttt{\tiny{[GM$|$GM]}}
                \texttt{continue} \\
            }
        }
    }
    & & \\ \\

    \theutterance \stepcounter{utterance}  
    & & & \multicolumn{2}{p{0.3\linewidth}}{
        \cellcolor[rgb]{0.9,0.9,0.9}{
            \makecell[{{p{\linewidth}}}]{
                \texttt{\tiny{[GM$|$GM]}}
                \texttt{I value the pepper the most, followed by the tomato, and the bags are worth the least to me. How do you value these items?} \\
            }
        }
    }
    & & \\ \\

    \theutterance \stepcounter{utterance}  
    & & \multicolumn{4}{p{0.6\linewidth}}{
        \cellcolor[rgb]{0.9,0.9,0.9}{
            \makecell[{{p{\linewidth}}}]{
                \texttt{\tiny{[GM$\rangle$P2]}}
                \texttt{You are playing a negotiation game in which you have to agree on how to divide a set of items among you and another player.} \\
\\ 
\texttt{Rules:} \\
\texttt{(a) You and the other player are given a set of items. Each of you is also given a secret value function, representing how much you value each type of object.} \\
\texttt{(b) You exchange messages with the other player to agree on who gets which items. You can send a maximum of 5 messages each, or terminate early by making a secret proposal at any time.} \\
\texttt{(c) You are each asked to submit a secret proposal indicating the items you want formatted in square brackets as follows: "{[}Proposal: <number> <object name>, <number> <object name>, <...>{]}"} \\
\texttt{(d) If your proposals are complementary, i.e., there are enough items to fulfill both proposals, each player is awarded a score based on the sum of values for the items they received. Otherwise, both of you get zero points.} \\
\texttt{(e) Your goal is to maximize the score you receive.} \\
\\ 
\texttt{Let us start.} \\
\\ 
\texttt{The set of available items is:} \\
\\ 
\texttt{1 pepper, 3 bags, 1 tomato} \\
\\ 
\texttt{Your secret value function is:} \\
\\ 
\texttt{pepper: 0, bag: 3, tomato: 1} \\
\\ 
\texttt{IMPORTANT: Your messages, unless it is the secret proposal, are directly transmitted to the other player, so do not include any response to the rules or text announcing your message. To make a secret proposal, use the indicated format. Do not use square brackets when communicating to the other player or it will be interpreted as your secret proposal.} \\
\\ 
\\ 
\texttt{I value the pepper the most, followed by the tomato, and the bags are worth the least to me. How do you value these items?} \\
            }
        }
    }
    & \\ \\

    \theutterance \stepcounter{utterance}  
    & & & \multicolumn{4}{p{0.6\linewidth}}{
        \cellcolor[rgb]{0.9,0.9,0.9}{
            \makecell[{{p{\linewidth}}}]{
                \texttt{\tiny{[GM$\langle$P2]}}
                \texttt{I value the bags the most, followed by the tomato, and the pepper is worth nothing to me. Since you want the pepper and I want the bags, we could easily split those. The tomato is the only item we both seem to want somewhat. How would you like to handle the tomato?} \\
            }
        }
    }
     \\ \\

    \theutterance \stepcounter{utterance}  
    & & & \multicolumn{2}{p{0.3\linewidth}}{
        \cellcolor[rgb]{0.9,0.9,0.9}{
            \makecell[{{p{\linewidth}}}]{
                \texttt{\tiny{[GM$|$GM]}}
                \texttt{continue} \\
            }
        }
    }
    & & \\ \\

    \theutterance \stepcounter{utterance}  
    & & & \multicolumn{2}{p{0.3\linewidth}}{
        \cellcolor[rgb]{0.9,0.9,0.9}{
            \makecell[{{p{\linewidth}}}]{
                \texttt{\tiny{[GM$|$GM]}}
                \texttt{I value the bags the most, followed by the tomato, and the pepper is worth nothing to me. Since you want the pepper and I want the bags, we could easily split those. The tomato is the only item we both seem to want somewhat. How would you like to handle the tomato?} \\
            }
        }
    }
    & & \\ \\

    \theutterance \stepcounter{utterance}  
    & & \multicolumn{4}{p{0.6\linewidth}}{
        \cellcolor[rgb]{0.9,0.9,0.9}{
            \makecell[{{p{\linewidth}}}]{
                \texttt{\tiny{[P1$\langle$GM]}}
                \texttt{I value the bags the most, followed by the tomato, and the pepper is worth nothing to me. Since you want the pepper and I want the bags, we could easily split those. The tomato is the only item we both seem to want somewhat. How would you like to handle the tomato?} \\
            }
        }
    }
    & \\ \\

    \theutterance \stepcounter{utterance}  
    & \multicolumn{4}{p{0.6\linewidth}}{
        \cellcolor[rgb]{0.9,0.9,0.9}{
            \makecell[{{p{\linewidth}}}]{
                \texttt{\tiny{[P1$\rangle$GM]}}
                \texttt{That sounds like a great plan. Since the pepper is worthless to you and the bags are low value to me, we can easily swap those. For the tomato, since we both value it, maybe we can flip a coin or you take it if you value it significantly more than I do? Or perhaps we can agree that you take the tomato and I take the pepper and one bag, while you take the other two bags? Let me know what works best for you.} \\
            }
        }
    }
    & & \\ \\

    \theutterance \stepcounter{utterance}  
    & & & \multicolumn{2}{p{0.3\linewidth}}{
        \cellcolor[rgb]{0.9,0.9,0.9}{
            \makecell[{{p{\linewidth}}}]{
                \texttt{\tiny{[GM$|$GM]}}
                \texttt{continue} \\
            }
        }
    }
    & & \\ \\

    \theutterance \stepcounter{utterance}  
    & & & \multicolumn{2}{p{0.3\linewidth}}{
        \cellcolor[rgb]{0.9,0.9,0.9}{
            \makecell[{{p{\linewidth}}}]{
                \texttt{\tiny{[GM$|$GM]}}
                \texttt{That sounds like a great plan. Since the pepper is worthless to you and the bags are low value to me, we can easily swap those. For the tomato, since we both value it, maybe we can flip a coin or you take it if you value it significantly more than I do? Or perhaps we can agree that you take the tomato and I take the pepper and one bag, while you take the other two bags? Let me know what works best for you.} \\
            }
        }
    }
    & & \\ \\

    \theutterance \stepcounter{utterance}  
    & & \multicolumn{4}{p{0.6\linewidth}}{
        \cellcolor[rgb]{0.9,0.9,0.9}{
            \makecell[{{p{\linewidth}}}]{
                \texttt{\tiny{[GM$\rangle$P2]}}
                \texttt{That sounds like a great plan. Since the pepper is worthless to you and the bags are low value to me, we can easily swap those. For the tomato, since we both value it, maybe we can flip a coin or you take it if you value it significantly more than I do? Or perhaps we can agree that you take the tomato and I take the pepper and one bag, while you take the other two bags? Let me know what works best for you.} \\
            }
        }
    }
    & \\ \\

    \theutterance \stepcounter{utterance}  
    & & & \multicolumn{4}{p{0.6\linewidth}}{
        \cellcolor[rgb]{0.9,0.9,0.9}{
            \makecell[{{p{\linewidth}}}]{
                \texttt{\tiny{[GM$\langle$P2]}}
                \texttt{That sounds like a fair approach. Since I value the bags the most (3 points each) and the tomato is worth 1 point to me, while the pepper is 0, I would prefer to keep all three bags and the tomato if possible, but I know you value the tomato too.} \\
\\ 
\texttt{However, looking at your suggestion: if I take the tomato and two bags, and you take the pepper and one bag, that works perfectly.} \\
\texttt{{-} You get: 1 pepper (0 value to me, but you value it), 1 bag (3 value to me, but you value it less).} \\
\texttt{{-} I get: 2 bags (6 value), 1 tomato (1 value). Total 7 for me.} \\
\\ 
\texttt{Wait, let's re{-}evaluate based on your values. You said you value pepper most, then tomato, then bags. I value bags most, then tomato, then pepper (0).} \\
\\ 
\texttt{If we do:} \\
\texttt{{-} You: 1 pepper, 1 tomato, 1 bag} \\
\texttt{{-} Me: 2 bags} \\
\\ 
\texttt{Or:} \\
\texttt{{-} You: 1 pepper, 1 bag} \\
\texttt{{-} Me: 2 bags, 1 tomato} \\
\\ 
\texttt{Actually, since there are 3 bags total, and I value them highly, I'd prefer to get as many as possible. But you need the pepper and tomato.} \\
\\ 
\texttt{Let's try this split:} \\
\texttt{{-} You get: 1 pepper, 1 tomato, 1 bag} \\
\texttt{{-} I get: 2 bags} \\
\\ 
\texttt{This way, you get your top two items (pepper and tomato) plus a bag, and I get two bags which are my highest value. Since the tomato is worth 1 to me and 1 to you (assuming your "followed by tomato" means it's your second highest, but you didn't specify exact numbers, just order), and bags are 3 to me and low to you, this seems efficient.} \\
\\ 
\texttt{Does that work for you? If so, I'll make a secret proposal.} \\
\\ 
\texttt{{[}Proposal: 2 bag, 1 tomato{]}} \\
            }
        }
    }
     \\ \\

    \theutterance \stepcounter{utterance}  
    & & & \multicolumn{2}{p{0.3\linewidth}}{
        \cellcolor[rgb]{0.9,0.9,0.9}{
            \makecell[{{p{\linewidth}}}]{
                \texttt{\tiny{[GM$|$GM]}}
                \texttt{proposal} \\
            }
        }
    }
    & & \\ \\

    \theutterance \stepcounter{utterance}  
    & & & \multicolumn{2}{p{0.3\linewidth}}{
        \cellcolor[rgb]{0.9,0.9,0.9}{
            \makecell[{{p{\linewidth}}}]{
                \texttt{\tiny{[GM$|$GM]}}
                [0, 2, 1]
            }
        }
    }
    & & \\ \\

    \theutterance \stepcounter{utterance}  
    & & \multicolumn{4}{p{0.6\linewidth}}{
        \cellcolor[rgb]{0.9,0.9,0.9}{
            \makecell[{{p{\linewidth}}}]{
                \texttt{\tiny{[P1$\langle$GM]}}
                \texttt{The other player has submitted a secret proposal. It is now your turn to submit a proposal of your own.} \\
            }
        }
    }
    & \\ \\

    \theutterance \stepcounter{utterance}  
    & \multicolumn{4}{p{0.6\linewidth}}{
        \cellcolor[rgb]{0.9,0.9,0.9}{
            \makecell[{{p{\linewidth}}}]{
                \texttt{\tiny{[P1$\rangle$GM]}}
                \texttt{{[}Proposal: 1 pepper, 1 tomato{]}} \\
            }
        }
    }
    & & \\ \\

    \theutterance \stepcounter{utterance}  
    & & & \multicolumn{2}{p{0.3\linewidth}}{
        \cellcolor[rgb]{0.9,0.9,0.9}{
            \makecell[{{p{\linewidth}}}]{
                \texttt{\tiny{[GM$|$GM]}}
                \texttt{proposal} \\
            }
        }
    }
    & & \\ \\

    \theutterance \stepcounter{utterance}  
    & & & \multicolumn{2}{p{0.3\linewidth}}{
        \cellcolor[rgb]{0.9,0.9,0.9}{
            \makecell[{{p{\linewidth}}}]{
                \texttt{\tiny{[GM$|$GM]}}
                [1, 0, 1]
            }
        }
    }
    & & \\ \\

    \theutterance \stepcounter{utterance}  
    & & & \multicolumn{2}{p{0.3\linewidth}}{
        \cellcolor[rgb]{0.9,0.9,0.9}{
            \makecell[{{p{\linewidth}}}]{
                \texttt{\tiny{[GM$|$GM]}}
                [[1, 0, 1], [0, 2, 1]]
            }
        }
    }
    & & \\ \\

    \theutterance \stepcounter{utterance}  
    & \multicolumn{4}{p{0.6\linewidth}}{
        \cellcolor[rgb]{0.95,0.95,0.95}{
            \makecell[{{p{\linewidth}}}]{
                \texttt{\tiny{[P2$\rightarrow$P1]}}
                \texttt{Post{-}episode Reflexion (not sent to GM):} \\
\texttt{New plan: I failed because I did not secure a complementary agreement before submitting my proposal. The other player explicitly stated they value bags the most and pepper the least, yet I proposed taking the pepper and the tomato, leaving them with only bags. While this might seem complementary, I failed to confirm the specific split of the three bags or ensure the other player's proposal matched my expectation of them taking all bags. My mistake was assuming the split without explicit confirmation of the bag distribution. Next time, I will explicitly propose a full division in my negotiation message (e.g., "I take the pepper and tomato; you take all 3 bags") and wait for their explicit agreement before submitting the secret proposal. If they submit a proposal first, I will carefully parse their request to ensure my proposal covers the remaining items exactly, rather than guessing. Specifically, I will ensure my proposal includes the items I want (pepper, tomato) and leaves the items they want (all 3 bags) for them, but only after confirming this specific allocation is acceptable to them.} \\
            }
        }
    }
    & & \\ \\

    \theutterance \stepcounter{utterance}  
    & & & \multicolumn{4}{p{0.6\linewidth}}{
        \cellcolor[rgb]{0.95,0.95,0.95}{
            \makecell[{{p{\linewidth}}}]{
                \texttt{\tiny{[P1$\rightarrow$P2]}}
                \texttt{Post{-}episode Reflexion (not sent to GM):} \\
\texttt{New plan: I failed because I misinterpreted the opponent's stated preferences and proposed a split that likely conflicted with their actual needs, leading to non{-}complementary proposals. Specifically, the opponent stated they value the pepper most, then the tomato, then bags. I incorrectly assumed they would accept a split where they got the pepper and tomato but I kept the tomato as well in my proposal logic, or I failed to secure the tomato for them when they clearly valued it highly. My proposal `{[}Proposal: 2 bag, 1 tomato{]}` left the opponent with only 1 pepper and 1 bag (since there are 3 bags total), denying them the tomato they valued second most.} \\
\\ 
\texttt{Next time, I will:} \\
\texttt{1.  **Strictly map the opponent's stated order to the available items:** If they say "Pepper > Tomato > Bags", they need the Pepper and Tomato to maximize their score.} \\
\texttt{2.  **Propose a split that gives the opponent their top priorities:** I should offer them the Pepper and the Tomato.} \\
\texttt{3.  **Keep the items I value most:** Since I value Bags (3) > Tomato (1) > Pepper (0), and the opponent wants the Tomato, I must negotiate to keep the Bags.} \\
\texttt{4.  **Formulate the correct proposal:** I will propose taking all 3 bags for myself (`{[}Proposal: 3 bag{]}`) and explicitly agree in the chat that the opponent takes the Pepper and Tomato. This ensures the proposals are complementary (3 bags + 1 pepper + 1 tomato = all items) and maximizes my score (9 points) while giving the opponent their high{-}value items, ensuring they accept the deal. I will not include the tomato in my proposal if the opponent values it highly, as that creates a conflict.} \\
            }
        }
    }
     \\ \\

\end{supertabular}
}

\end{document}
